\exnumber{03}
\extitle{UX I - Ex 4 - Parcours utilisateur}
\exfiles{ex03/fichier Figma, nomme Name\_FirstName\_DeliverableName\_Date\_VersionNumber}
\exauthorize{Ceux nommés dans l'exercice}
\makeheaderfiles

\textbf{Instructions :}

À partir de votre carte persona et des scripts d'entretien, concevez le user journey de Pierre.

\textbf{Livrables :}

\begin{itemize}
\item Un fichier Figma avec le parcours utilisateur de Pierre.
\item Déposez le fichier Figma sur Github, nommé ainsi : « Name\_FirstName\_DeliverableName\_Date\_VersionNumber ». Exemple : « Doe\_John\_MidFiPrototype\_05242024\_V1 ».
\item Ajoutez un titre dans le fichier.
\item Soignez l'aspect global du fichier. Il doit être clair, simple et lisible. Vous pouvez regarder canva.com, sans surcharger le dessin de détails.
\end{itemize}

\textbf{Pistes :}

\begin{itemize}
\item Vous pouvez utiliser Figma.
\item Mélangez des insights et des détails tirés des scripts d'entretien pour rédiger un parcours cohérent avec la personnalité de Pierre.
\item Vous pouvez, par exemple, décrire le premier marathon de Pierre. La première phase pourrait être la préparation, la deuxième la course elle-même, et la troisième comment il se sent par rapport à la course une fois qu'il l'a terminée.
\item Découpez l'histoire étape par étape et définissez les actions de Pierre.
\item Concentrez-vous sur chaque action et définissez les pensées de Pierre ainsi que son émotion principale parmi la liste.
\item Adaptez la ligne aux pensées et émotions de Pierre.
\item Définir les opportunités d'améliorer l'expérience globale de Pierre chaque fois que vous identifiez un point de douleur.
\end{itemize}

\textbf{Ressources :}

La plupart des articles que vous trouverez porteront sur les user journeys ou customer journeys. Gardez à l'esprit que vous pouvez adapter cette méthodologie à n'importe quel persona dans n'importe quel type de situation ou de processus, même à l'école, en dehors de l'école, dans le monde numérique, etc.

\begin{itemize}
\item Outil : Template User Journey Mapping sur Figma \url{https://www.figma.com/file/sXdUDoTXwItU7dyQYVkY1w/Tool---User-Journey-Map}
\item Une Customer Journey Map Lean UX \url{https://uxdesign.cc/proto-journey-a-lean-ux-customer-journey-map-30ea3a241edc}
\item Comment créer une Customer Journey Map \url{https://www.youtube.com/watch?v=mSxpVRo3BLg}
\item Journey Mapping : comment le construire \url{https://www.lucidchart.com/blog/how-to-build-customer-journey-maps}
\item Journey Mapping : Template sur Figma \url{https://www.figma.com/file/5tittjiznRCWTNR2xI5FIT/Untitled?node-id=0%3A1}
\item Exemple \url{https://d2slcw3kip6qmk.cloudfront.net/marketing/blog/2017Q3/SEO-initiative-customer-journey-mapping/CustomerJourneyMap1.png}
\end{itemize}

\textbf{Citation :}

« Il s'agit de surprendre les clients en action, et de leur fournir une information hautement pertinente et contextualisée » Paul Maritz | PDG Pivotal

